% !TEX program = pdflatex
% Rapport complet — Pipeline Medallion Météo Maroc (Parties 1 à 6)
% Module : Architecture des Données — 4IASDG2 — EMSI — 2024/2025
% Compiler : pdflatex rapport_complet_4IASDG2.tex (2 fois)

\documentclass[12pt,a4paper]{report}

\usepackage[T1]{fontenc}
\usepackage[utf8]{inputenc}
\usepackage[french]{babel}
\usepackage[top=2.5cm,bottom=2.5cm,left=2.5cm,right=2.5cm,headheight=15pt]{geometry}
\usepackage{setspace}
\onehalfspacing
\usepackage{graphicx}
\usepackage{float}
\usepackage{booktabs}
\usepackage{array}
\usepackage{longtable}
\usepackage{xcolor}
\usepackage{listings}
\usepackage{hyperref}
\usepackage{fancyhdr}
\usepackage{titlesec}
\usepackage{enumitem}
\usepackage{caption}
\usepackage{tikz}
\usetikzlibrary{arrows.meta, positioning, shapes.geometric, fit, backgrounds}

\definecolor{emsiBlue}{RGB}{0,82,155}
\definecolor{bronzeColor}{RGB}{205,127,50}
\definecolor{silverColor}{RGB}{160,160,160}
\definecolor{goldColor}{RGB}{218,165,32}
\definecolor{codeBg}{RGB}{245,245,245}

\hypersetup{colorlinks=true,linkcolor=emsiBlue,urlcolor=emsiBlue,citecolor=emsiBlue}

\pagestyle{fancy}
\fancyhf{}
\fancyhead[L]{\small Architecture des Données — 4IASDG2}
\fancyhead[R]{\small EMSI — 2024/2025}
\fancyfoot[C]{\thepage}

\lstset{
    backgroundcolor=\color{codeBg}, frame=single, basicstyle=\ttfamily\small,
    breaklines=true, numbers=left, numberstyle=\tiny\color{gray},
    keywordstyle=\color{emsiBlue}\bfseries
}

\titleformat{\chapter}[display]
  {\normalfont\huge\bfseries\color{emsiBlue}}
  {\chaptertitlename\ \thechapter}{20pt}{\Huge}

% ─────────────────────────────────────────────────────────────
\begin{document}

% ═══ PAGE DE GARDE ═══
\begin{titlepage}
\centering
\vspace*{0.8cm}
{\Large\textbf{ÉCOLE MAROCAINE DES SCIENCES DE L'INGÉNIEUR}}\\[0.2cm]
{\large EMSI — Campus Rabat}\\[1.8cm]
\rule{\linewidth}{1.5pt}\\[0.4cm]
{\Huge\bfseries\color{emsiBlue} Rapport de Projet}\\[0.3cm]
{\LARGE\bfseries Pipeline de Données Météorologiques}\\[0.15cm]
{\LARGE\bfseries Architecture Medallion — Bronze à Power BI}\\[0.4cm]
\rule{\linewidth}{1.5pt}\\[1.5cm]

\begin{tabular}{rl}
\textbf{Module :} & Architecture des Données \\[0.4cm]
\textbf{Filière :} & 4IASD — Groupe 4IASDG2 \\[0.4cm]
\textbf{Année :} & 2024--2025 \\[0.4cm]
\textbf{Étudiant(s) :} & \rule{7.5cm}{0.4pt} \\[0.5cm]
\textbf{Encadrant :} & \rule{7.5cm}{0.4pt} \\[0.5cm]
\textbf{Date :} & \rule{7.5cm}{0.4pt} \\
\end{tabular}

\vfill
{\large\textbf{Flux global :}}\\[0.2cm]
{\small
API Weather $\rightarrow$ Airbyte $\rightarrow$ Bronze $\rightarrow$ Silver $\rightarrow$ Gold $\rightarrow$ Airflow $\rightarrow$ Power BI
}
\vspace{1cm}
\end{titlepage}

% ═══ RÉSUMÉ ═══
\chapter*{Résumé}
\addcontentsline{toc}{chapter}{Résumé}

Ce rapport documente la conception, l'implémentation et la validation d'un \textbf{pipeline de données complet} pour l'analyse météorologique de \textbf{huit villes marocaines}. Le projet suit l'\textbf{architecture Medallion} en six phases successives : ingestion Bronze via \textbf{Airbyte} et \textbf{MinIO}, transformation Silver en Python, enrichissement Gold avec calcul de KPI, tests de qualité automatisés, orchestration quotidienne par \textbf{Apache Airflow}, et visualisation décisionnelle dans \textbf{Power BI}.

L'infrastructure repose sur \textbf{Docker Desktop}, un stockage objet \textbf{MinIO} compatible S3, et un déploiement local d'Airbyte via \textbf{abctl}. Les résultats incluent un score météo sur 100 points, un classement des villes, un dashboard interactif et un pipeline automatisé exécuté chaque jour à 6h00.

\textbf{Mots-clés :} Data Lake, Medallion, Bronze, Silver, Gold, Airbyte, MinIO, Airflow, Power BI, OpenWeatherMap, KPI, 4IASDG2.

\tableofcontents
\listoffigures
\listoftables
\newpage

% ═══ CHAPITRE 1 — INTRODUCTION ═══
\chapter{Introduction}

\section{Contexte}

Les systèmes d'information décisionnels modernes s'appuient sur des pipelines de données structurés, capables de collecter, transformer, contrôler et exposer des données fiables. Dans le cadre du module \textbf{Architecture des Données} (4IASDG2), ce projet vise à mettre en œuvre un tel pipeline appliqué aux \textbf{données météorologiques} de villes marocaines.

\section{Problématique}

Comment construire un pipeline de bout en bout — de l'ingestion brute à la visualisation décisionnelle — en respectant les bonnes pratiques d'ingénierie des données (Medallion, qualité, orchestration) ?

\section{Objectifs}

\begin{enumerate}[leftmargin=1.5cm]
    \item \textbf{Partie 1 — Bronze :} Ingérer des données météo brutes depuis une API Weather vers MinIO.
    \item \textbf{Partie 2 — Silver :} Nettoyer et structurer les données JSON brutes.
    \item \textbf{Partie 3 — Gold :} Calculer des KPI métier (score météo, classement).
    \item \textbf{Partie 4 — Qualité :} Valider automatiquement chaque couche.
    \item \textbf{Partie 5 — Airflow :} Automatiser le pipeline (cron 6h00).
    \item \textbf{Partie 6 — Power BI :} Créer un dashboard décisionnel.
\end{enumerate}

\section{Périmètre}

\begin{table}[H]
\centering
\caption{Villes marocaines du pipeline}
\begin{tabular}{@{}ll@{}}
\toprule
\textbf{Ville} & \textbf{Usage} \\
\midrule
Casablanca, Rabat, Marrakech, Fès & Grandes agglomérations \\
Tanger, Agadir, Meknès, Oujda & Couverture géographique nationale \\
\bottomrule
\end{tabular}
\end{table}

% ═══ CHAPITRE 2 — ARCHITECTURE ═══
\chapter{Architecture et environnement technique}

\section{Architecture Medallion}

\begin{center}
\begin{tikzpicture}[
    layer/.style={draw, rounded corners, minimum width=2.8cm, minimum height=0.9cm, align=center, font=\small},
    arr/.style={-{Stealth}, thick}
]
    \node[layer, fill=bronzeColor!25] (b) {Bronze\\{\tiny Raw JSON}};
    \node[layer, fill=silverColor!25, right=0.6cm of b] (s) {Silver\\{\tiny Clean JSON}};
    \node[layer, fill=goldColor!25, right=0.6cm of s] (g) {Gold\\{\tiny Analytics}};
    \node[layer, fill=emsiBlue!15, right=0.6cm of g] (bi) {Power BI\\{\tiny Dashboard}};
    \draw[arr] (b) -- (s); \draw[arr] (s) -- (g); \draw[arr] (g) -- (bi);
    \node[above=0.8cm of b, font=\small] {Airbyte + MinIO};
    \node[above=0.8cm of s, font=\small] {Python / Pandas};
    \node[above=0.8cm of g, font=\small] {Python KPI};
    \node[above=0.8cm of bi, font=\small] {JSON Export};
\end{tikzpicture}
\end{center}

\section{Stack technologique}

\begin{table}[H]
\centering
\caption{Composants du pipeline}
\begin{tabular}{@{}p{3.5cm}p{4cm}p{6cm}@{}}
\toprule
\textbf{Rôle} & \textbf{Technologie} & \textbf{Version / Détail} \\
\midrule
Conteneurisation & Docker Desktop & 29.x \\
Ingestion & Airbyte (abctl + kind) & 2.1.0 \\
Stockage objet & MinIO & Ports 9000 / 9001 \\
Transformation & Python, Pandas, Boto3 & 3.11+ \\
Orchestration & Apache Airflow & 2.8.1 \\
Visualisation & Power BI Desktop & Dernière version \\
Source API & OpenWeatherMap & Plan gratuit \\
\bottomrule
\end{tabular}
\end{table}

\section{Chemins de stockage MinIO}

\begin{table}[H]
\centering
\caption{Organisation des données dans MinIO}
\begin{tabular}{@{}ll@{}}
\toprule
\textbf{Couche} & \textbf{Chemin} \\
\midrule
Bronze & \texttt{bronze/weather\_raw/YYYY-MM-DD/<ville>.json} \\
Silver & \texttt{silver/weather\_clean/YYYY-MM-DD/weather\_clean.json} \\
Gold   & \texttt{gold/weather\_analytics/YYYY-MM-DD/weather\_analytics.json} \\
\bottomrule
\end{tabular}
\end{table}

% ═══ CHAPITRE 3 — PARTIE 1 BRONZE ═══
\chapter{Partie 1 — Ingestion Bronze (Airbyte + MinIO)}

\section{Objectif}

Collecter les données météorologiques brutes depuis l'API OpenWeatherMap et les déposer \textbf{sans transformation} dans MinIO.

\section{Mise en œuvre}

\subsection{Déploiement MinIO}

\begin{lstlisting}[language=bash]
docker run -d --name minio -p 9000:9000 -p 9001:9001
  -e MINIO_ROOT_USER=minioadmin
  -e MINIO_ROOT_PASSWORD=minioadmin
  minio/minio server /data --console-address ":9001"
\end{lstlisting}

\subsection{Déploiement Airbyte}

\begin{lstlisting}[language=bash]
abctl.exe local install --low-resource-mode --no-browser
abctl.exe local credentials
# Interface : http://localhost:8000
\end{lstlisting}

\subsection{Configuration Airbyte}

\textbf{Source OpenWeather :} clé API, latitude/longitude par ville, unités métriques.

\textbf{Destination S3/MinIO :}
\begin{itemize}[leftmargin=1.5cm]
    \item Endpoint : \url{http://host.docker.internal:9000}
    \item Bucket : \texttt{bronze}
    \item Prefix : \texttt{weather\_raw/}
    \item Credentials : minioadmin / minioadmin
\end{itemize}

\section{Résultat attendu}

Huit fichiers JSON bruts par jour, un par ville, contenant température, humidité, vent, pression, description météo et horodatages.

% ═══ CHAPITRE 4 — PARTIE 2 SILVER ═══
\chapter{Partie 2 — Transformation Silver}

\section{Objectif}

Transformer les données brutes en données \textbf{propres, structurées et exploitables}.

\section{Travail réalisé}

Le script \texttt{silver\_transform.py} effectue les opérations suivantes :

\begin{enumerate}[leftmargin=1.5cm]
    \item Lecture des JSON Bronze depuis MinIO (Boto3) ;
    \item Extraction des champs : ville, température, humidité, vent, description ;
    \item Conversion des types (\texttt{float}, \texttt{int}, timestamp ISO) ;
    \item Suppression des valeurs nulles et incohérentes ;
    \item Conservation de \texttt{event\_timestamp} et \texttt{ingestion\_timestamp}.
\end{enumerate}

\section{Schéma Silver}

\begin{table}[H]
\centering
\caption{Structure du fichier \texttt{weather\_clean.json}}
\begin{tabular}{@{}lll@{}}
\toprule
\textbf{Champ} & \textbf{Type} & \textbf{Description} \\
\midrule
ville & string & Nom de la ville \\
temperature & float & Température en °C \\
humidite & int & Humidité relative (\%) \\
vent\_vitesse & float & Vitesse du vent (m/s) \\
event\_timestamp & datetime & Date de l'observation météo \\
ingestion\_timestamp & datetime & Date d'ingestion dans le pipeline \\
\bottomrule
\end{tabular}
\end{table}

\section{Stockage}

Fichier unique agrégé : \texttt{silver/weather\_clean/YYYY-MM-DD/weather\_clean.json} contenant 8 enregistrements (un par ville).

% ═══ CHAPITRE 5 — PARTIE 3 GOLD ═══
\chapter{Partie 3 — Enrichissement Gold}

\section{Objectif}

Créer des \textbf{données métier} pour la prise de décision : score météo, classement, meilleure ville.

\section{Calcul du score météo (KPI)}

Le score sur 100 points combine trois critères :

\begin{table}[H]
\centering
\caption{Grille de notation du score météo}
\begin{tabular}{@{}p{4cm}p{2cm}p{7cm}@{}}
\toprule
\textbf{Critère} & \textbf{Points max} & \textbf{Condition optimale} \\
\midrule
Température & 40 & Entre 20 et 28 °C \\
Humidité    & 30 & Inférieure à 50 \% \\
Vent        & 30 & Inférieur à 5 m/s \\
\bottomrule
\end{tabular}
\end{table}

\section{Classement et meilleure ville}

Les villes sont triées par score décroissant. En cas d'ex-aequo, le classement partagé est appliqué (ex. deux villes à 90/100 = rang 1). La \textbf{meilleure ville du jour} est celle avec le rang minimal.

\section{Exemple de résultats}

\begin{table}[H]
\centering
\caption{Classement météo — exemple}
\begin{tabular}{@{}lrrrr@{}}
\toprule
\textbf{Ville} & \textbf{Temp. (°C)} & \textbf{Humidité (\%)} & \textbf{Score} & \textbf{Rang} \\
\midrule
Tanger    & 21.64 & 46 & 90 & 1 \\
Agadir    & 22.25 & 41 & 90 & 1 \\
Marrakech & 19.04 & 55 & 75 & 3 \\
Oujda     & 19.35 & 49 & 65 & 4 \\
Rabat     & 22.04 & 60 & 65 & 4 \\
Casablanca& 17.19 & 88 & 60 & 6 \\
Fès       & 17.14 & 72 & 60 & 6 \\
Meknès    & 18.49 & 63 & 60 & 6 \\
\bottomrule
\end{tabular}
\end{table}

\section{Stockage Gold}

\texttt{gold/weather\_analytics/YYYY-MM-DD/weather\_analytics.json} + export local \texttt{powerbi/weather\_analytics.json} pour Power BI.

% ═══ CHAPITRE 6 — PARTIE 4 QUALITÉ ═══
\chapter{Partie 4 — Qualité des données}

\section{Objectif}

Garantir la \textbf{fiabilité} des données à chaque étape du pipeline.

\section{Tests implémentés}

\begin{table}[H]
\centering
\caption{Matrice des tests qualité}
\begin{tabular}{@{}p{2cm}p{5cm}p{6cm}@{}}
\toprule
\textbf{Couche} & \textbf{Test} & \textbf{Critère de succès} \\
\midrule
Bronze & Existence fichiers & 8 fichiers JSON présents \\
Silver & Types et colonnes  & Champs obligatoires, types corrects, humidité 0--100 \\
Gold   & KPI et classement  & Score $\in$ [0,100], 8 villes, rangs cohérents \\
\bottomrule
\end{tabular}
\end{table}

\section{Exécution}

\begin{lstlisting}[language=bash]
python scripts/data_quality.py --layer bronze
python scripts/data_quality.py --layer silver
python scripts/data_quality.py --layer gold
\end{lstlisting}

En cas d'échec, le script lève une exception et bloque le pipeline (fail-fast).

% ═══ CHAPITRE 7 — PARTIE 5 AIRFLOW ═══
\chapter{Partie 5 — Orchestration Airflow}

\section{Objectif}

\textbf{Automatiser} l'exécution complète du pipeline sans intervention manuelle.

\section{DAG \texttt{weather\_morocco\_pipeline}}

\begin{center}
\begin{tikzpicture}[
    task/.style={draw, rounded corners, fill=emsiBlue!10, minimum width=2.2cm, minimum height=0.7cm, font=\scriptsize, align=center},
    arr/.style={-{Stealth}, thick}
]
    \node[task] (a) {ingestion\\bronze};
    \node[task, right=0.4cm of a] (b) {test\\bronze};
    \node[task, right=0.4cm of b] (c) {transform\\silver};
    \node[task, right=0.4cm of c] (d) {test\\silver};
    \node[task, right=0.4cm of d] (e) {enrich\\gold};
    \node[task, right=0.4cm of e] (f) {test\\gold};
    \draw[arr] (a)--(b)--(c)--(d)--(e)--(f);
\end{tikzpicture}
\end{center}

\section{Paramètres du DAG}

\begin{table}[H]
\centering
\caption{Configuration Airflow}
\begin{tabular}{@{}ll@{}}
\toprule
\textbf{Paramètre} & \textbf{Valeur} \\
\midrule
DAG ID & \texttt{weather\_morocco\_pipeline} \\
Schedule & \texttt{0 6 * * *} (quotidien à 6h00) \\
Executor & SequentialExecutor \\
Image Docker & \texttt{apache/airflow:2.8.1} \\
Interface & \url{http://localhost:8080} \\
Tags & etl, maroc, weather \\
\bottomrule
\end{tabular}
\end{table}

% ═══ CHAPITRE 8 — PARTIE 6 POWER BI ═══
\chapter{Partie 6 — Reporting Power BI}

\section{Objectif}

Créer un \textbf{dashboard décisionnel} connecté à la couche Gold.

\section{Connexion aux données}

\begin{enumerate}[leftmargin=1.5cm]
    \item Power BI Desktop $\rightarrow$ Obtenir des données $\rightarrow$ JSON ;
    \item Fichier : \texttt{powerbi/weather\_analytics.json} ;
    \item Charger la table des 8 villes.
\end{enumerate}

\section{Visuels créés}

\begin{table}[H]
\centering
\caption{Dashboard Power BI — 4 visuels}
\begin{tabular}{@{}p{3.5cm}p{4cm}p{6cm}@{}}
\toprule
\textbf{Visuel} & \textbf{Type} & \textbf{Champs} \\
\midrule
Température par ville & Graphique barres & ville, temperature \\
Classement & Table & ville, score\_meteo, classement \\
Meilleure ville & Carte KPI & ville (rang = 1) \\
Carte Maroc & Carte géographique & ville, score\_meteo \\
\bottomrule
\end{tabular}
\end{table}

% ═══ CHAPITRE 9 — DIFFICULTÉS ═══
\chapter{Difficultés rencontrées}

\begin{enumerate}[leftmargin=1.5cm]
    \item \textbf{Déploiement Airbyte :} échec initial avec docker-compose ; résolu via abctl.
    \item \textbf{Espace disque :} installation Helm longue ($>$25 min) ; images volumineuses.
    \item \textbf{Secret manquant :} \texttt{airbyte-auth-secrets not found} — installation Helm incomplète ; relance de \texttt{local install}.
    \item \textbf{Endpoint MinIO :} utilisation de \texttt{host.docker.internal:9000} pour Airbyte dans Docker.
\end{enumerate}

% ═══ CHAPITRE 10 — CONCLUSION ═══
\chapter{Conclusion}

Ce projet a permis de mettre en place un \textbf{pipeline de données complet} conforme à l'énoncé 4IASDG2 (parties 1 à 6). L'architecture Medallion garantit la séparation des responsabilités entre ingestion, transformation, enrichissement et reporting.

Les apports principaux :
\begin{itemize}[leftmargin=1.5cm]
    \item Maîtrise de l'ingestion ELT avec Airbyte et MinIO ;
    \item Implémentation Python des couches Silver et Gold ;
    \item Automatisation et contrôle qualité via Airflow ;
    \item Visualisation décisionnelle dans Power BI.
\end{itemize}

\textbf{Perspectives :} ajout de prévisions météo, alertes automatiques, migration vers le cloud (AWS S3, Azure Data Lake), et intégration Spark pour le traitement à grande échelle.

% ═══ ANNEXES ═══
\appendix

\chapter{Structure du projet}

\begin{lstlisting}
weather_project/
  config/          # settings.py, minio_client.py
  scripts/         # bronze, silver, gold, data_quality
  dags/            # weather_morocco_pipeline.py
  powerbi/         # weather_analytics.json, dashboard.pbix
  docker-compose.yml
  requirements.txt
\end{lstlisting}

\chapter{Commandes de relance}

\begin{lstlisting}[language=bash]
docker start minio
abctl.exe local install --low-resource-mode --no-browser
python scripts/run_pipeline.py
docker compose up -d   # Airflow
# Power BI : actualiser le fichier JSON
\end{lstlisting}

\chapter{Liens utiles}

\begin{tabular}{@{}p{4cm}p{10cm}@{}}
Docker & \url{https://www.docker.com/products/docker-desktop/} \\
Airbyte & \url{https://docs.airbyte.com} \\
MinIO & \url{https://min.io/docs/minio/container/index.html} \\
OpenWeather & \url{https://openweathermap.org/api} \\
Airflow & \url{https://airflow.apache.org/docs/} \\
Power BI & \url{https://powerbi.microsoft.com} \\
\end{tabular}

\begin{thebibliography}{9}
\addcontentsline{toc}{chapter}{Bibliographie}
\bibitem{medallion} Databricks, \textit{Medallion Architecture}, 2023.
\bibitem{airbyte} Airbyte Inc., \textit{Documentation OSS}, 2024.
\bibitem{minio} MinIO Inc., \textit{Object Storage Docs}, 2024.
\bibitem{airflow} Apache Software Foundation, \textit{Airflow Documentation}, 2024.
\end{thebibliography}

\end{document}
